\begin{lemma}[3.4.1.2]
\label{IV.2.3.4.1.2-fr}
Sean $A$ un anillo local noetheriano de dimensión $>0$, $\mathfrak p_i$ los
ideales primos mínimos de $A$, $\mathfrak m$ su ideal maximal, $t$ un
elemento de $\mathfrak m$ tal que $A/tA$ sea artiniano. Entonces, para todo
$A$-módulo de tipo finito $M$, se tiene
\end{lemma}

\begin{env}[3.4.1.3]
\label{IV.2.3.4.1.3-fr}
\[
\operatorname{long}(M/tM)
\geqslant\sum_i\operatorname{long}(M_{\mathfrak p_i})~;
\]
\end{env}

\textit{además, para que los dos miembros de} (3.4.1.3) \textit{sean iguales, es necesario y
suficiente que se satisfagan las condiciones siguientes~:}

\oldpage[IV-2]{47}

\begin{enumerate}
\item[$\alpha$)] \textit{$t$ es $M$-regular~;}
\item[$\beta$)] \textit{para todo $i$ tal que $M_{\mathfrak p_i}\neq0$, la imagen de $t$
en $A/\mathfrak p_i$ genera el ideal maximal de este anillo} (lo que implica
que $A/\mathfrak p_i$ es un \emph{anillo de valoración discreta}).
\end{enumerate}

\begin{proof}
Como $A$ no es de dimensión $0$ y $A/tA$ es artiniano, necesariamente se tiene
$\dim(A)=1$ $(\mathbf{0},16.3.4)$ y $t\notin\mathfrak p_i$ para todo
$i$~: el ideal principal $(t)$ es, por tanto, un ideal de definición de $A$, y
contiene por consiguiente una potencia de su ideal maximal $\mathfrak m$. Sea $N$
el submódulo de los elementos de $M$ anulados por una potencia de $t$ (o por
una potencia de $\mathfrak m$, lo que, como se acaba de ver, equivale a lo mismo)~;
si se pone $P=M/N$, $t$ es $P$-\emph{regular}, pues la relación $tx\in N$
para un $x\in M$ implica $t^k(tx)=0$ para un entero $k$, y por tanto $x\in N$.
Hecho esto, se tiene el

\begin{lemma}[3.4.1.4]
\label{IV.2.3.4.1.4-fr}
Sean $A$ un anillo,
\[
0\longrightarrow M'\longrightarrow M\longrightarrow M''\longrightarrow0
\]
una sucesión exacta de $A$-módulos. Si $t\in A$ es $M''$-regular, la sucesión
\[
0\longrightarrow M'/tM'\longrightarrow M/tM
\longrightarrow M''/tM''\longrightarrow0
\]
es exacta.
\end{lemma}

\begin{proof}
Como $M/tM=M\otimes_A(A/tA)$, basta demostrar la exactitud en $M'/tM'$~;
ahora bien, si la imagen $x\in M$ de un elemento $x'$ de $M'$ es tal que
$x=ty$ con $y\in M$, se deduce, para las imágenes $x''$, $y''$ de $x$, $y$ en
$M''$ que $x''=ty''$~; pero como $x''=0$, la hipótesis implica $y''=0$, de
modo que $y$ es la imagen de un elemento $y'\in M'$, y la relación $x=ty$
implica $x'=ty'$ puesto que $M'\to M$ es inyectiva.
\end{proof}

Una vez establecido este lema, se obtiene la relación
\end{proof}

\begin{env}[3.4.1.5]
\label{IV.2.3.4.1.5-fr}
\[
\operatorname{long}(M/tM)
=\operatorname{long}(N/tN)+\operatorname{long}(P/tP).
\]
\end{env}

